\chapter{Freight, Weather and Other Corners}\label{ch:m3:freight-weather-and-other-corners}

A gas utility in Chicago earns its money in winter. A cold winter fills its pipes and its
accounts; a warm one leaves it with gas bought for customers who did not need it. In the
winter of 2011/12 the city's airport recorded 3{,}932 \omterm{def:m3:freight-weather-and-other-corners:dd}{heating degree days}, the fewest in
thirty-six winters; two years later it recorded 5{,}999, the most. The utility can buy a
contract that pays it for each degree day by which the winter falls short of a strike,
settled on the temperatures measured at that airport. Freight and weather are markets in
things that are not commodities at all: the use of a ship for a voyage, the temperature of
a season. This chapter covers how ships are chartered and how freight is hedged, how
weather is turned into a traded index and priced, and what it takes for any such corner to
become a market.

\section{Freight: charters and routes}

Bulk commodities travel in ships hired for the purpose. The hire takes two forms.

\begin{definition}[Voyage charter, time charter]\label{def:m3:freight-weather-and-other-corners:charter}
A \emph{voyage charter}\index{voyage charter} hires a ship to carry a cargo between named ports
for a freight rate per tonne; the shipowner pays the voyage costs (fuel, port charges). A
\emph{time charter}\index{time charter} hires a ship for a period, for a daily rate of hire; the
charterer directs its voyages and pays their costs.
\end{definition}

\begin{definition}[Time-charter equivalent]\label{def:m3:freight-weather-and-other-corners:tce}
The \emph{time-charter equivalent}\index{time-charter equivalent} (TCE) of a voyage is its freight
revenue less its voyage costs, divided by its duration in days: the daily hire at which a \omterm{def:m3:freight-weather-and-other-corners:charter}{time
charter} would earn the owner the same.
\end{definition}

\begin{example}[From a voyage to a daily rate]\label{ex:m3:freight-weather-and-other-corners:tce}
A bulk carrier carries 170{,}000 tonnes of iron ore at \$12.00 a tonne: freight of \$2.04 million.
Fuel costs \$600{,}000 and port charges \$150{,}000, and the round voyage takes 45 days. Its TCE is
$(2{,}040{,}000 - 750{,}000)/45 = \$28{,}667$ a day.
\end{example}

\begin{omfigure}
\centering
\begin{tikzpicture}[font=\small,
  pty/.style={draw,thick,rounded corners=2pt,align=center,minimum width=2.6cm,minimum height=1.1cm},
  arr/.style={-{Stealth[length=2.2mm]},thick}]
\node[pty,draw=omDef,fill=omDef!8] (own) at (0,0) {shipowner};
\node[pty,draw=omThm,fill=omThm!8] (chr) at (9.6,0) {charterer};
\draw[arr] ([yshift=4pt]chr.west) -- node[above=2pt,font=\footnotesize] {voyage charter: \$ per tonne} ([yshift=4pt]own.east);
\draw[arr] ([yshift=-4pt]chr.west) -- node[below=2pt,font=\footnotesize] {time charter: \$ per day} ([yshift=-4pt]own.east);
\node[font=\footnotesize,align=center,anchor=north] at (0,-0.9) {pays fuel and ports\\under a voyage charter};
\node[font=\footnotesize,align=center,anchor=north] at (9.6,-0.9) {pays fuel and ports\\under a time charter};
\end{tikzpicture}

\omcaption{The two ways of hiring a ship, and who bears the voyage costs. The \omterm{def:m3:freight-weather-and-other-corners:tce}{time-charter
equivalent} converts one into the other. Schematic.}\label{fig:m3:freight-weather-and-other-corners:charters}
\end{omfigure}

\begin{definition}[Route index]\label{def:m3:freight-weather-and-other-corners:route}
A \emph{route index}\index{route index} is a daily assessment, by a panel of shipbrokers under an
index provider's rules, of the rate for a standard ship on a standard route, in dollars per tonne or
per day; indices for several routes are averaged into a basket (for Capesize bulk carriers, a
weighted average of five time-charter routes).
\end{definition}

\begin{dated}{2026-09}{The freight index provider}\label{dat:m3:freight-weather-and-other-corners:baltic}
The Baltic Exchange, which publishes the main dry and wet freight indices, was acquired by Singapore
Exchange; the acquisition was completed in November 2016.
\end{dated}

\section{Forward freight agreements}

\begin{definition}[Forward freight agreement]\label{def:m3:freight-weather-and-other-corners:ffa}
A \emph{forward freight agreement}\index{forward freight agreement} (FFA) is a cash-settled forward on
a \omterm{def:m3:freight-weather-and-other-corners:route}{route index} for a future period, usually a month: at settlement the buyer receives, per day (or
tonne) of the agreed quantity, the average of the index over the period minus the fixed rate agreed.
Most are cleared.
\end{definition}

A shipowner that expects to let its ship in six months sells FFAs to lock its hire; a charterer or
a trading house that will need ships buys them. Because FFAs settle on a monthly average, they hedge
the average rate of the month, not the rate of the day a ship is fixed: the difference is \omterm{def:m3:physical-commodity-markets:basisrisk}{basis risk}
(\cref{ch:m3:physical-commodity-markets}).

\begin{omfigure}
\begin{tikzpicture}
\begin{axis}[width=11cm,height=4.4cm,
  xlabel={index day of the month},ylabel={thousand \$ per day},
  tick label style={font=\small},label style={font=\small},
  xmin=0,xmax=26,ymin=20,ymax=27,grid=major,grid style={gray!25},
  legend columns=3,legend style={font=\footnotesize,at={(0.5,-0.33)},anchor=north,draw=none,fill=none,/tikz/every even column/.append style={column sep=8pt}},
  table/col sep=comma]
\addplot[omDef,very thick,mark=*,mark size=1pt] table[x=day,y=index]{figdata/markets-3/11-freight-weather-and-other-corners/ffa.csv};
\addplot[omThm,thick,dashed,domain=0:26] {23.4};
\addplot[omProp,thick,dotted,domain=0:26] {21.0};
\legend{daily index,monthly average (settlement),fixed rate of the FFA}
\end{axis}
\end{tikzpicture}

\omcaption{A time-charter FFA's settlement month, illustrative: the index rises from \$21{,}000 to
\$25{,}800 a day over 25 index days, averages \$23{,}400, and the buyer of 30 days at \$21{,}000
receives \$72{,}000. Data: the chapter's tutorial.}\label{fig:m3:freight-weather-and-other-corners:ffa}
\end{omfigure}

\begin{dated}{2026-09}{Freight derivatives}\label{dat:m3:freight-weather-and-other-corners:ffa}
Capesize freight futures and FFAs cleared on Singapore Exchange settle in cash on the arithmetic
average of the Baltic Exchange's daily assessments of the five-route time-charter basket in the
contract month; the same indices are cleared as futures by other exchanges.
\end{dated}

\section{Weather: degree days and weather derivatives}

Energy demand follows temperature, and temperature can be measured without dispute at a named
station. The industry's index turns a season's temperatures into one number.

\begin{definition}[Heating degree day, cooling degree day]\label{def:m3:freight-weather-and-other-corners:dd}
With $\theta_i$ the average of a day's maximum and minimum temperatures in degrees Fahrenheit, the
day's \emph{heating degree days}\index{heating degree day} are $\mathrm{HDD}_i = \max(0, 65 - \theta_i)$
and its \emph{cooling degree days}\index{cooling degree day} $\mathrm{CDD}_i = \max(0, \theta_i - 65)$;
a period's index is the sum over its days.
\end{definition}

\begin{definition}[Weather derivative]\label{def:m3:freight-weather-and-other-corners:wd}
A \emph{weather derivative}\index{weather derivative} is a contract whose payoff is a function of a
weather index at a named station over a period: a swap that pays a tick per degree day above or below
a strike, or an option that pays only on one side, usually with a cap.
\end{definition}

\begin{dated}{2026-09}{Exchange-traded weather}\label{dat:m3:freight-weather-and-other-corners:cme}
\omterm{def:m3:freight-weather-and-other-corners:wd}{Weather derivatives} traded over the counter from 1997; the first standalone transaction is described
as a Koch Energy and Enron HDD swap for a Milwaukee winter. The Chicago Mercantile Exchange listed HDD
futures in September 1999 and CDD futures in January 2000, on indices with a 65°F base. Its rulebook lists thirteen
US cities, each read at a named airport weather station (Chicago at O'Hare), sums a calendar month's degree days, and
sets the contract at USD~20 per index point with a one-point tick.
\end{dated}

\begin{omfigure}
\begin{tikzpicture}
\begin{axis}[width=11cm,height=4.8cm,ybar,bar width=4pt,
  xlabel={winter (year of the January)},ylabel={HDD, November--March},
  tick label style={font=\small},label style={font=\small},
  xmin=1990,xmax=2027,ymin=3000,ymax=6500,grid=major,grid style={gray!25},
  xtick={1991,1996,2001,2006,2011,2016,2021,2026},xticklabel style={/pgf/number format/1000 sep={}},
  table/col sep=comma]
\addplot[fill=omDef!55,draw=omDef] table[x=winter,y=hdd]{figdata/markets-3/11-freight-weather-and-other-corners/hdd.csv};
\draw[omThm,thick,dashed] (axis cs:1990,4911.8) -- (axis cs:2027,4911.8);
\node[omThm,font=\footnotesize,anchor=south west] at (axis cs:1990.5,6000) {mean 4{,}912};
\end{axis}
\end{tikzpicture}

\omcaption{\omterm{def:m3:freight-weather-and-other-corners:dd}{Heating degree days} at Chicago O'Hare, November to March, winters 1990/91 to 2025/26 (daily
average of maximum and minimum, 65°F base). The warmest winter was 2011/12, the coldest 2013/14; the
fitted trend is $-125$ HDD a decade. Data: NOAA GHCN-Daily, station USW00094846.}\label{fig:m3:freight-weather-and-other-corners:hdd}
\end{omfigure}

\section{Pricing weather: burn analysis}

A degree-day index has no underlying asset to hedge with and no forward price to replicate: weather
cannot be stored or traded. Pricing therefore starts from history.

\begin{definition}[Burn analysis]\label{def:m3:freight-weather-and-other-corners:burn}
\emph{Burn analysis}\index{burn analysis} prices a weather contract by applying its payoff to the index
values of past seasons, as if each had happened again, and taking the average payout as the fair value,
with the distribution of payouts as its risk.
\end{definition}

\begin{method}[Burn analysis with a trend]\label{met:m3:freight-weather-and-other-corners:burn}
(i) Compute the index for each of the last $n$ complete seasons at the contract's station.
(ii) Fit a linear trend and move each season onto the trend's value for the season priced (climate and
station changes make old winters look colder than the next one will be). (iii) Apply the payoff to each
adjusted season. (iv) The fair value is the mean payout; quote a margin above it for the risk, measured
by the frequency and the high quantiles of the payouts.
\end{method}

\begin{example}[The utility's put]\label{ex:m3:freight-weather-and-other-corners:put}
On the last thirty winters at O'Hare, detrended to the winter of 2026/27, the burn estimate of the
season's HDD is 4{,}677. A put struck 10\% lower, at 4{,}209, paying \$20{,}000 per HDD with a
\$20 million cap, would have paid in 20\% of the adjusted winters: \$0.61 million on average, \$1.60
million in the 90th-percentile winter and \$9.30 million in the warmest. Without the detrending the
estimate is 4{,}871 and the average payout \$0.56 million.
\end{example}

\begin{omfigure}
\begin{tikzpicture}
\begin{axis}[width=11cm,height=4.4cm,ybar,bar width=5pt,
  xlabel={winter (detrended to 2026/27)},ylabel={payout, \$ million},
  tick label style={font=\small},label style={font=\small},
  xmin=1996,xmax=2027,ymin=0,ymax=10,grid=major,grid style={gray!25},
  xtick={1997,2002,2007,2012,2017,2022,2026},xticklabel style={/pgf/number format/1000 sep={}},
  table/col sep=comma]
\addplot[fill=omThm!55,draw=omThm] table[x=winter,y=payout]{figdata/markets-3/11-freight-weather-and-other-corners/burn.csv};
\end{axis}
\end{tikzpicture}

\omcaption{\omterm{def:m3:freight-weather-and-other-corners:burn}{Burn analysis} of the utility's put: the payout it would have received in each of the last
thirty winters, adjusted to the trend's 2026/27 level. Six winters pay; 2011/12 pays the most. Data:
the chapter's tutorial, from NOAA GHCN-Daily.}\label{fig:m3:freight-weather-and-other-corners:burn}
\end{omfigure}

\omterm{def:m3:freight-weather-and-other-corners:burn}{Burn analysis} is crude: thirty seasons say little about a one-in-fifty winter, and the trend is itself
uncertain. Desks refine it by simulating daily temperatures and by blending in the season's forecast
when it is available, but the price remains an insurance price, set by the capital of those who write
the risk.

\section{What makes a market tradable at all}

Freight and weather are corners because only a few conditions made them markets, and markets in other
corners have failed for lack of one. The durable ones share five features:
\begin{enumerate}
\item \emph{A hedging need on both sides}: shipowners and charterers, utilities and energy traders who
are warm-winter losers and winners.
\item \emph{An index nobody can move}: a panel assessment under published rules, a thermometer at an
airport.
\item \emph{Standard contracts}: routes, stations, periods and ticks that concentrate liquidity.
\item \emph{Clearing}: margin that lets strangers trade without credit lines.
\item \emph{Speculators who will warehouse the risk}: without them hedgers must find an exact opposite,
which rarely exists.
\end{enumerate}
A contract that lacks the second fails when the index can be influenced; one that lacks the first or the
fifth never trades.

\section{Tutorial: degree days and a burn price}\label{tut:m3:freight-weather-and-other-corners:burn}

\begin{tutorial}
\textbf{Goal.} From daily station temperatures, build winter HDD totals, detrend them, and price a
degree-day put by \omterm{def:m3:freight-weather-and-other-corners:burn}{burn analysis}. \textbf{End state:}
\cref{fig:m3:freight-weather-and-other-corners:hdd,fig:m3:freight-weather-and-other-corners:burn} and the
numbers of \cref{ex:m3:freight-weather-and-other-corners:put}.
\begin{tutsteps}
\item \textbf{Indices and payoffs.}
\omcode{code/firm/degreeday/firm_degreeday.py}{21}{46}{Degree days, swap and put payoffs, and burn analysis.}\label{lst:m3:freight-weather-and-other-corners:burn}
\item \textbf{The trend.}
\omcode{code/firm/degreeday/firm_degreeday.py}{49}{56}{Detrending past seasons to the season priced.}\label{lst:m3:freight-weather-and-other-corners:detrend}
\item \textbf{Run} \texttt{m3\_weather.stats()}, \texttt{utility\_put()} and \texttt{fig\_weather.py}.
\end{tutsteps}
\textbf{What to change next.} Price the same put on the last ten winters only, and see how much the price
moves; price a CDD swap for July at the same station.
\end{tutorial}

\section{Build: degree days and freight settlement}\label{bld:m3:freight-weather-and-other-corners:degreeday}

\begin{build}
\bfield{Purpose} The miniature firm writes weather protection for utilities and hedges its own freight:
it needs station indices, burn prices and the settlement of freight agreements.
\bfield{Interface} \texttt{c10\_to\_f}, \texttt{daily\_average}, \texttt{hdd}, \texttt{cdd};
\texttt{swap\_payoff}, \texttt{put\_payoff(index, strike, tick, cap)}; \texttt{burn(history, payoff)};
\texttt{detrend(history, years, target\_year)}; \texttt{ffa\_settlement(daily\_index, fixed\_rate, days,
lots)}.
\bfield{Rules} Temperatures in °F, base 65; incomplete seasons are excluded, never filled; the burn
quantile is the empirical one; FFA settlement uses the arithmetic mean of published index days.
\bfield{Acceptance tests} \texttt{code/firm/degreeday/tests/}: unit conversions and degree days; payoffs
with a cap; a burn price on four seasons; detrending a linear series; an FFA settlement.
\bfield{Stretch} Daily temperature simulation (seasonal mean plus autoregressive noise); a portfolio of
stations with correlated winters; FFA curves with seasonality.
\end{build}

\begin{omsources}
\item NOAA National Centers for Environmental Information, GHCN-Daily, station USW00094846 (public domain).
\item CME Group, ``Overview of weather markets'' and rulebook chapter 403 (degree days index futures).
\item Baltic Exchange, news of 2016 on the acquisition by SGX; SGX freight contract specifications.
\item S.~Jewson and A.~Brix, \emph{Weather Derivative Valuation}, Cambridge University Press, 2005.
\item CME Rulebook, chapter 403, CME Degree Days Index Futures (Internet Archive copy).
\end{omsources}

\section{Exercises}

\begin{exercise}[$\star$]\label{exo:m3:freight-weather-and-other-corners:1}
A day has a maximum of 40°F and a minimum of 22°F. Give its HDD and CDD.
\end{exercise}

\begin{exercise}[$\star$]\label{exo:m3:freight-weather-and-other-corners:2}
A voyage earns \$1.5 million of freight, costs \$450{,}000 in fuel and \$100{,}000 in port charges, and takes
38 days. Give its TCE.
\end{exercise}

\begin{exercise}[$\star$]\label{exo:m3:freight-weather-and-other-corners:3}
Why does an FFA settle on a monthly average rather than on one day's index?
\end{exercise}

\begin{exercise}[$\star\star$]\label{exo:m3:freight-weather-and-other-corners:4}
An HDD swap has a strike of 4{,}700 and a tick of \$20{,}000. The winter comes in at 4{,}054. Who pays whom,
and how much?
\end{exercise}

\begin{exercise}[$\star\star$]\label{exo:m3:freight-weather-and-other-corners:5}
An owner sells 30 days of a time-charter FFA at \$21{,}000 a day; the month averages \$19{,}500. What does it
receive, and what has it hedged?
\end{exercise}

\begin{exercise}[$\star\star$]\label{exo:m3:freight-weather-and-other-corners:6}
Why is \omterm{def:m3:freight-weather-and-other-corners:burn}{burn analysis} on detrended winters more expensive, for this put, than on raw winters?
\end{exercise}

\begin{exercise}[$\star\star\star$]\label{exo:m3:freight-weather-and-other-corners:7}
\emph{Coding.} With \texttt{winters}, give the number of complete winters, their mean and standard deviation,
the warmest and coldest, and the fitted trend per decade.
\end{exercise}

\begin{exercise}[$\star\star\star$]\label{exo:m3:freight-weather-and-other-corners:8}
\emph{Find the flaw.} ``We price weather options with Black--Scholes on the HDD index, using its historical
volatility.''
\end{exercise}

\section{Problem: The Warm Winter}

\begin{problem}[{Weekend problem --- a Chicago utility's winter protection}]\label{pb:m3:freight-weather-and-other-corners:1}
A gas utility loses about \$20{,}000 of margin for each HDD by which a winter falls short of normal. It
considers an HDD put on O'Hare for November 2026 to March 2027, \$20{,}000 per HDD, capped at \$20 million.

\medskip\noindent\textbf{Part I --- The history.}
\begin{enumerate}
\item How many complete winters does the station provide since 1990/91, and what is their mean?
\item What is the standard deviation of the winter totals?
\item Which winters were the warmest and the coldest, with their totals?
\item What trend does a linear fit show?
\item Why should the price use a detrended history?
\end{enumerate}

\medskip\noindent\textbf{Part II --- The burn price.}
\begin{enumerate}[resume]
\item Give the detrended burn estimate of the 2026/27 winter.
\item Give the strike 10\% below it.
\item In what share of the adjusted winters does the put pay?
\item Give its average payout (the burn price) and its 90th-percentile payout.
\item What was its largest payout, and in which winter?
\end{enumerate}

\medskip\noindent\textbf{Part III --- The deal.}
\begin{enumerate}[resume]
\item Why will a seller ask more than the burn price?
\item What does the cap protect, and what does it cost the utility?
\item How does the utility's loss differ from the put's payout (basis)?
\item Would a swap be better than a put for the utility? For its shareholders?
\item Who might write the put, and how would they hedge it?
\end{enumerate}

\medskip\noindent\textbf{Part IV --- Judgement.}
\begin{enumerate}[resume]
\item What does thirty years of data leave out?
\item How would a long-range forecast change the price in October?
\item Why do weather markets stay small?
\item State the \emph{named result}: the burn-analysis strike and the one-in-ten-year payout of the put.
\item In one sentence: what is a \omterm{def:m3:freight-weather-and-other-corners:wd}{weather derivative}?
\end{enumerate}
\end{problem}

\section{Interview questions}

\begin{interviewq}[$\star$ \iqroles{trader}]\label{iq:m3:freight-weather-and-other-corners:1}
What is a \omterm{def:m3:freight-weather-and-other-corners:dd}{heating degree day}, and who buys HDD protection?
\end{interviewq}

\begin{interviewq}[$\star$ \iqroles{trader, researcher}]\label{iq:m3:freight-weather-and-other-corners:2}
What is a \omterm{def:m3:freight-weather-and-other-corners:ffa}{forward freight agreement}, and what does it settle on?
\end{interviewq}

\begin{interviewq}[$\star\star$ \iqroles{researcher}]\label{iq:m3:freight-weather-and-other-corners:3}
How would you price a weather option, and why not with Black--Scholes?
\end{interviewq}

\begin{interviewq}[$\star\star$ \iqroles{trader}]\label{iq:m3:freight-weather-and-other-corners:4}
A shipowner will have a ship available in six months. How does it hedge, and what risks remain?
\end{interviewq}

\begin{interviewq}[$\star\star$ \iqroles{risk}]\label{iq:m3:freight-weather-and-other-corners:5}
You sell HDD puts on ten US cities. What is your risk, and how do you measure it?
\end{interviewq}

\begin{interviewq}[$\star\star\star$ \iqroles{researcher, developer}]\label{iq:m3:freight-weather-and-other-corners:6}
Design a pipeline that prices degree-day contracts for any station daily, from raw weather data to quotes.
\end{interviewq}
